\section{Limitations}
\label{sec:limitations}

\paragraph{Scale.} Every experiment runs on one consumer GPU (RTX 5070, 12\,GB). To fit a five-hour
training budget we cut the pre-registered grid, which called for five seeds per cell and a full maze
grid, down to two seeds per method, three single-seed ablations, and a smaller gate-tuning budget for
PonderNet and the v1 gate than was pre-registered (\texttt{PREREGISTRATION.md}, deviations). With two
seeds the accuracy intervals are little more than the seed range, and the pre-registered Wilcoxon tests
cannot reach significance. Some comparisons that we report as ties may separate with more seeds, and
small differences should not be read as effects.

\paragraph{Mazes.} We could not train a maze backbone that passes Gate G1 within the budget; when
stopped, it reached 3\% exact match on training-size mazes. We therefore make no claims about stopping
rules on mazes, and whether local completion heads help when information travels along corridors
remains open.

\paragraph{Tasks.} The exact theory covers one-dimensional tasks whose determining radius has a closed
form. Addition and segmented parity are deliberately simple, so that $\LB(x)$ is known for every
instance, and they are not representative of natural-language reasoning. Our conclusions concern
recurrent models with a known, local growth rate. They do not transfer directly to models with global
attention, whose light cone covers the whole input after a single step.

\paragraph{Backbone dependence.} A stopping rule can only be as good as the trajectory it reads. On
segmented parity the backbone fails Gate G1, since its stable-oracle accuracy out of distribution is
below the pre-registered 95\%, so the parity results measure stopping on an imperfect backbone. The
progressive loss was also seed-dependent on addition. Termination detection and depth extrapolation
are entangled: a rule that looks premature may in fact be reacting to a backbone that never resolves.

\paragraph{What the theory does not say.} \Cref{thm:lb} rules out \emph{certified} early answers, not
correct ones. A rule that halts before $\LB(x)$ can be right on the data distribution, and on addition
the convergence detectors often are, because a carry wave keeps the hidden state moving until it lands.
This is why we report the premature-stop and premature-error rates of \cref{eq:premature} separately,
and we do not claim that premature halting is always harmful.

\paragraph{Selection and thresholds.} Thresholds are tuned on ID-val by taking the cheapest setting
within 0.5 points of the best. A different tolerance would trade accuracy against steps differently.
The gate hyperparameters (the ACT ponder weight $\tauact$, the PonderNet prior $\lambda_p$ and the v1 gate
weight $\lambda$) were tuned on a small grid with one seed, and stronger tuning of the jointly trained
gates could narrow some of the gaps we report.

\paragraph{Broader impact.} This is basic research on small synthetic tasks. Better termination
detection lowers inference compute, and we see no direct path to harm beyond that of adaptive-compute
models in general.
